\section*{Capítulo \ref{ch:b2:flow-balances} --- Balanços de quantidade de movimento e de energia em escoamentos}

\begin{solution}{exo:b2:flow-balances:1}
$s = \qty{7.1}{cm^2}$, $v = 0.025/7.1 \times 10^{-4} = \qty{35}{m/s}$; $F = D_mv = 25
\times 35 = \qty{880}{N}$ no esguicho e na parede; parede que recua: $\rho s(v -
u)^2 = 1000 \times 7.1 \times 10^{-4} \times 25.4^2 = \qty{460}{N}$.
\end{solution}

\begin{solution}{exo:b2:flow-balances:2}
$F = 250 \times 3000 = \qty{750}{kN}$; $a_0 = 750/60 - 9.81 = \qty{2.7}{m/s^2}$; a
\qty{20}{t}: $37.5 - 9.8 = \qty{28}{m/s^2}$. $\Delta v = 3000\ln3 = \qty{3.3}{km/s}$ no
espaço livre; menos $g\tau = \qty{1.6}{km/s}$: \qty{1.7}{km/s}.
\end{solution}

\begin{solution}{exo:b2:flow-balances:3}
$S = \qty{177}{cm^2}$: $PS = \qty{7.1}{kN}$, $\rho Sv^2 = \qty{71}{N}$. Cada componente
da força sobre o cotovelo vale $PS + \rho Sv^2 = \qty{7.1}{kN}$: \qty{10}{kN} ao longo
da bissetriz externa, $99\%$ disso de pressão. Com a saída aberta (pressão efetiva
nula ali): $7.1$ kN na direção da entrada e \qty{71}{N} na direção da
saída: \qty{7.1}{kN}, quase segundo o eixo de entrada.
\end{solution}

\begin{solution}{exo:b2:flow-balances:4}
$F = 80 \times 550 = \qty{44}{kN}$. Em voo: $F = 80(550 - 240) = \qty{25}{kN}$, $Fv =
\qty{6.0}{MW}$; $\mathcal P_{\text{cin}} = \tfrac12 \times 80 \times (550^2 - 240^2) =
\qty{9.8}{MW}$; $\eta_p = 0.61 = 2 \times 240/790$.
\end{solution}

\begin{solution}{exo:b2:flow-balances:5}
$D_m = \qty{500}{kg/s}$. (a) $F = 1000(60 - u)$ N, $\mathcal P = 1000u(60 - u)$; $u =
\qty{30}{m/s}$, $\omega = 37.5$ rad/s $= \qty{360}{rpm}$. (b) \qty{900}{kW} $=
\tfrac12D_mv^2$: $100\%$. (c) $F = D_m(v - u)(1 + \cos15^\circ) = \qty{29.5}{kN}$, $\mathcal P
= \qty{885}{kW}$, $98\%$. (d) $\Gamma = FR = \qty{24}{kN.m}$; com a roda parada a força
dobra ($\Gamma = \qty{48}{kN.m}$); em $u = v$ força e torque se anulam.
\end{solution}

\begin{solution}{exo:b2:flow-balances:6}
(a) $s = \qty{7.1}{mm^2}$ por bico, \qty{0.10}{L/s} em cada: $w = \qty{14}{m/s}$. (b)
$\Gamma_0 = D_mRw = 0.2 \times 0.15 \times 14 = \qty{0.42}{N.m}$. (c) $D_mR(w - R\omega) =
5 \times 10^{-3}$: $w - R\omega = \qty{0.17}{m/s}$, $\omega = \qty{93}{rad/s} \approx \qty{890}{rpm}$;
sem atrito, $R\omega = w$: \qty{900}{rpm}. (d) \qty{0.17}{m/s} para trás, e depois
zero: a água cai direto no chão.
\end{solution}

\begin{solution}{exo:b2:flow-balances:7}
(a) $D_m = \rho Av_i$, \omterm{thm:b2:flow-balances:momentum}{fluxo de quantidade de movimento} $D_m\cdot2v_i$: $T = 2\rho Av_i^2$; $A = \qty{78.5}{m^2}$,
$T = \qty{19.6}{kN}$: $v_i = \qty{10}{m/s}$. (b) $\mathcal P = \tfrac12D_m(2v_i)^2 = Tv_i =
\qty{200}{kW}$. (c) $\mathcal P = T^{3/2}/\sqrt{2\rho A}$: dobrar a área
divide a potência induzida por $\sqrt2$ --- um rotor maior move mais ar
mais devagar. (d) $A = \qty{0.20}{m^2}$, $v_i = \sqrt{9.81/0.47} = \qty{4.6}{m/s}$, $\mathcal P
= \qty{45}{W}$: \qty{45}{W/kg} contra \qty{100}{W/kg} --- a carga do disco $T/A$
é menor.
\end{solution}

\begin{solution}{exo:b2:flow-balances:8}
(a) $v_2 = \qty{1.0}{m/s}$; por Bernoulli, $\tfrac12\rho(16 - 1) = \qty{7.5}{kPa}$; de fato,
$\rho v_2(v_1 - v_2) = \qty{3.0}{kPa}$; perda de \qty{4.5}{kPa} $= \qty{0.46}{m}$. (b) $6$ kPa
recuperados, \qty{1.5}{kPa} perdidos. (c) $\tfrac12\rho v_1^2 = \qty{8}{kPa}$, \qty{0.82}{m}. (d)
$D_V = \qty{7.9}{L/s}$: \qty{35}{W} e \qty{63}{W}.
\end{solution}

\begin{solution}{exo:b2:flow-balances:9}
(a) $v = 0.01/2.83 \times 10^{-3} = \qty{3.5}{m/s}$, $v^2/2g = \qty{0.64}{m}$; atrito,
$0.025 \times 1000 \times 0.64 = \qty{16}{m}$, conexões \qty{1.3}{m}: \qty{17.3}{m}. (b)
$H = 30 + 17.3 + 0.64 = \qty{48}{m}$; $\rho gD_VH = \qty{4.7}{kW}$; \qty{7.2}{kW} no
eixo. (c) Entrada (sem perdas na sucção): $P_0 - \rho g(5 + 0.64) = \qty{45}{kPa}$;
saída, $45 + 470 = \qty{515}{kPa}$. (d) $P_{\text{ent}} \ge \qty{2.3}{kPa}$: $h \le (100 -
2.3)/9.81 - 0.64 = \qty{9.3}{m}$ (menos com perdas na sucção; $7$--\qty{8}{m} na
prática).
\end{solution}

\begin{solution}{exo:b2:flow-balances:10}
(a) $v = \sqrt{2 \times 4 \times 10^5/1000} = \qty{28}{m/s}$; $s = \qty{3.8}{cm^2}$, $D_m =
\qty{10.8}{kg/s}$, $F = D_mv = 2\Delta P\,s = \qty{300}{N}$; $a = 300/1.1 - 9.8 =
\qty{260}{m/s^2}$. (b) Ar a \qty{1.5}{L}: $P = \qty{3.3}{bar}$, $v = \sqrt{2 \times 2.3 \times
10^5/1000} = \qty{22}{m/s}$. (c) Celeridade média de saída $\approx \qty{22}{m/s}$: propulsão $\approx 10^{-3}
/(3.8 \times 10^{-4} \times 22) = \qty{0.12}{s}$; a \omterm{prop:b2:flow-balances:pelton}{equação do foguete} com $v_e \approx \qty{22}{m/s}$ dá
$\Delta v \approx 22\ln(1.1/0.1) = \qty{50}{m/s}$. (d) Os \qty{2.5}{bar} de ar restantes
sopram para fora e acrescentam um pouco; água demais deixa pouco ar e a
pressão desaba cedo, água de menos deixa pouca massa a lançar ---
cerca de um terço do volume é o ideal.
\end{solution}

\begin{solution}{exo:b2:flow-balances:11}
(a) $\Gamma = D_mr_1v_{\theta1} = 3 \times 10^4 \times 20 = \qty{600}{kN.m}$, $\omega = 15.7$ rad/s,
$\mathcal P = \qty{9.4}{MW}$. (b) $H = \mathcal P/\rho gD_V = \qty{32}{m}$. (c) $\Gamma = 3 \times
10^4(20 - 1.5) = \qty{555}{kN.m}$, $\mathcal P = \qty{11.6}{MW}$ (carga de \qty{39}{m}); o
giro de saída leva $v_{\theta2}^2/2g = \qty{0.46}{m}$, cerca de $1\%$. (d) O jato de uma Pelton
está à pressão atmosférica: precisa de grande carga para ser rápido, e sua
vazão é limitada pelo bocal; um rotor Francis fica cheio de água sob
pressão e passa grandes vazões com cargas moderadas.
\end{solution}

\begin{solution}{exo:b2:flow-balances:12}
(a) $v_1h_1 = v_2h_2 = q$. (b) Força hidrostática por unidade de largura numa
seção: $\int_0^h\rho g(h - z)\dd z = \tfrac12\rho gh^2$; quantidade de movimento: $\tfrac12\rho g(h_1^2
- h_2^2) = \rho q(v_2 - v_1)$; divida por $\rho$ e use $q = v_1h_1 = v_2h_2$.
(c) Com $v_2 = v_1h_1/h_2$ e $h_1 \ne h_2$: $\tfrac12g(h_1 + h_2)h_2 = v_1^2h_1$,
ou seja, $(h_2/h_1)^2 + h_2/h_1 - 2\mathrm{Fr}_1^2 = 0$; $h_2 > h_1$ exige $\mathrm{Fr}_1 > 1$.
(d) $\Delta H = (h_1 + v_1^2/2g) - (h_2 + v_2^2/2g) = (h_2 - h_1)^3/4h_1h_2$.
$\mathrm{Fr}_1 = 2.86$, $h_2/h_1 = 3.57$, $h_2 = \qty{0.71}{m}$, $v_2 = \qty{1.1}{m/s}$,
$\Delta H = 0.514^3/0.57 = \qty{0.24}{m}$; $\rho gq\,\Delta H = \qty{1.9}{kW}$ por metro de
largura.
\end{solution}

\begin{solution}{pb:b2:flow-balances:1}
\textbf{1.} $F = (D_a + D_f)v_e = 51 \times 600 = \qty{30.6}{kN}$.

\textbf{2.} $F = 51 \times 600 - 50 \times 250 = \qty{18.1}{kN}$; $Fv = \qty{4.5}{MW}$.

\textbf{3.} $\mathcal P_{\text{cin}} = \tfrac12 \times 51 \times 600^2 - \tfrac12 \times 50 \times 250^2 =
\qty{7.6}{MW}$; $\eta_p = 0.59$. Com $D_f = 0$: $\eta_p = 2v(v_e - v)/(v_e^2 - v^2)
= 2v/(v_e + v) = 500/850$.

\textbf{4.} Calor de \qty{43}{MW}: térmico $7.6/43 = 18\%$; global $4.5/43 = 10\%$.

\textbf{5.} $\tfrac12 \times 500(v_e'^2 - 250^2) = 7.6 \times 10^6$: $v_e' = \qty{305}{m/s}$;
$F = 500 \times 55 = \qty{27.5}{kN}$; $\eta_p = 500/555 = 0.90$. Mover mais ar mais
devagar dá mais empuxo para a mesma potência: daí os ventiladores enormes.

\textbf{6.} Ele precisa do ar externo tanto como fluido de trabalho quanto como
oxidante; um foguete leva os dois.

\textbf{7.} O ar entra a \qty{70}{m/s} (para trás, no referencial do motor)
e sai com uma componente para a frente $150\cos45^\circ = \qty{106}{m/s}$: a
quantidade de movimento dada ao ar aponta para a frente, $500(106 + 70) = \qty{88}{kN}$ de
frenagem; o mesmo fluxo expelido para trás daria $500(150 - 70) =
\qty{40}{kN}$ de empuxo.

\textbf{8.} Em $\dd t$ o sistema (foguete $m$ + gás $D_m\dd t$ ejetado a
$v - v_e$): $(m - D_m\dd t)(v + \dd v) + D_m\dd t(v - v_e) - mv = -mg\,\dd t$,
donde $m\,\dd v/\dd t = D_mv_e - mg$.

\textbf{9.} $D_m = \qty{1333}{kg/s}$; $F = \qty{4.0}{MN}$; $a = 13.3 - 9.8 =
\qty{3.5}{m/s^2}$ na decolagem, $40 - 9.8 = \qty{30}{m/s^2}$ no fim da queima.

\textbf{10.} $v = v_e\ln(m_0/m) - gt$: $3000\ln3 - 9.81 \times 150 = 3296 - 1472 =
\qty{1.8}{km/s}$; a perda por gravidade é de \qty{1.5}{km/s}.

\textbf{11.} Estágio 1: $3000\ln(300/200) = \qty{1.2}{km/s}$; descarta \qty{20}{t}; estágio
2: $3000\ln(180/80) = \qty{2.4}{km/s}$: \qty{3.6}{km/s} contra \qty{3.3}{km/s}
do estágio único (com gravidade a diferença aumenta).

\textbf{12.} $v_e = 7800/\ln3 = \qty{7.1}{km/s}$: nenhum propelente químico
chega lá (hidrogênio--oxigênio dá \qty{4.5}{km/s}); o escalonamento é obrigatório.

\textbf{13.} O fluxo vertical de quantidade de movimento do jato $D_mv = \qty{4.0}{MN}$ é
retirado e um fluxo horizontal igual é criado: uma força de \qty{4}{MN} para baixo
mais \qty{4}{MN} de lado, \qty{5.7}{MN} ao todo.

\textbf{14.} O balanço no motor inclui $(P_e - P_0)S_e$ na
seção de saída; como $P_0$ cai com a altitude, esse termo cresce: o empuxo
sobe uns $10$--$15\%$ no vácuo.

\textbf{15.} $D_m = \qty{3e4}{kg/s}$ por roda: $F = D_m(v - u)(1 + \cos15^\circ) =
1.97 \times 3 \times 10^4(62.6 - u)$; $\mathcal P = Fu$, máxima em $u = v/2 = \qty{31.3}{m/s}$.

\textbf{16.} $\omega = u/R = 20.9$ rad/s $= \qty{199}{rpm}$; $F = \qty{1.85}{MN}$, $\Gamma = FR
= \qty{2.8}{MN.m}$.

\textbf{17.} $f = pn/60$: $p = 3000/199 \approx 15$ pares de polos, $n = \qty{200}{rpm}$.

\textbf{18.} $\mathcal P = Fu = \qty{58}{MW}$; potência do jato $\tfrac12D_mv^2 = \qty{59}{MW}$:
$98\%$; $2 \times 58 \times 0.95 = \qty{110}{MW}$.

\textbf{19.} Na entrada: $Rv = \qty{94}{m^2/s}$ por kg. Na saída: celeridade azimutal absoluta
$u - (v - u)\cos15^\circ = \qty{1.1}{m/s}$, ou seja, \qty{1.6}{m^2/s}. $\Gamma = D_m(94 -
1.6) = \qty{2.8}{MN.m}$ --- como em 16.

\textbf{20.} Quando $u \to v$ a celeridade relativa, a força e o torque
se anulam: nada freia a roda, que dispara; o defletor
tira o jato das conchas em menos de um segundo, enquanto a válvula-agulha
fecha devagar o bastante para evitar o golpe de aríete.

\textbf{21.} $v = \qty{5.7}{m/s}$, $v^2/2g = \qty{1.6}{m}$; perda $0.015 \times 133 \times 1.6 =
\qty{3.3}{m}$; carga da bomba $200 + 3.3 + 1.6 = \qty{205}{m}$.

\textbf{22.} $\rho gD_VH = 9810 \times 40 \times 205 = \qty{80}{MW}$; \qty{91}{MW} elétricos.

\textbf{23.} $P = P_0 + \rho g(200 + 3.3) + \tfrac12\rho v^2 \approx \qty{21}{bar}$ absolutos,
contra \qty{20.3}{bar} em modo turbina: agora as perdas se somam à
carga estática em vez de subtrair-se dela.

\textbf{24.} $V = 40 \times 28800 = \qty{1.15e6}{m^3}$, $E = \rho gV \times 200 = \qty{2.3e12}{J}
= \qty{630}{MWh}$; consumidos $91 \times 8 = \qty{730}{MWh}$; recuperados $630 \times 0.9 \times
0.95 = \qty{540}{MWh}$: ida e volta, $74\%$.

\textbf{25.} Quantidade de movimento: o empuxo do motor a jato e o do foguete, a
força na Pelton. Momento angular: o torque da turbina (Euler). Energia: a
carga e a potência da bomba, o rendimento do armazenamento. Massa: toda vazão,
e a perda de massa do foguete.
\end{solution}
